\section{Application: legislative incidence at the district level}
\label{sec:application}

The pipeline's output is in production use. The 57{,}240-record sparse
support selected in \autoref{sp:sec:results} ships, after
post-selection refit, as the certified United States release of
\populace{}, pinned by immutable release tag and consumed through
\policyenginepy{}'s managed data bundle alongside a pinned
model version. A companion paper \citep{ghenis2026obbba} uses it to
decompose the One Big Beautiful Bill Act of 2025 --- eighteen tax
provisions and three program-participation scenarios --- against a
TCJA-expiration counterfactual, reporting incidence at the national,
state, and congressional-district level, with poverty and inequality
statistics computed on the same file.

Precision about that lineage matters for replication. The shipped
support is the production $L_0$ selection carried forward across
builds by identity join --- including a nine-household swap that
protects Keogh-distribution carriers (\populace{} issue \#434) ---
and its publication weights are refit, gates off, on the certified
national-and-state target surface (5{,}659 targets) for 6{,}000
epochs, reaching a 1.77\% calibration loss. The 32{,}633-target probe
of \autoref{sp:sec:results}, which adds the congressional-district
families, is the research configuration on the same candidate
universe. On the shipped weights the maximum household weight is
50{,}665 and the effective sample size 9{,}174; 40{,}600 of the
57{,}240 records carry weight below one household, so the calibrator,
not the record budget, determines the effective support. Release
manifests record the selection lineage, target-registry version, and
engine pins.

Three aspects of that deployment bear directly on this paper's
claims. First, the imputed inputs stage one exists to provide are
load-bearing: the companion analysis's tip, overtime, auto-loan, and
program-participation provisions run entirely on imputed or
enhanced columns, and its health-inclusive resource measure depends
on the modeled program dollars the enhanced file supports.

Second, the value of target-informed selection is directly
observable in the district layer. The companion paper's district
estimates use weights calibrated against roughly 24{,}000
district-level targets. Recomputing its district averages on
assignment-only weights --- the same national file, with records
assigned to districts but with no district targets in the
calibration --- produces a minimum district cell of 34 records, a
district with a negative average net-income change, and an
interdecile district ratio of 3.0 against the calibrated layer's
2.2: tail artifacts of sparse, unbalanced support rather than
incidence. Random assignment is to district estimation what random
subsampling is to national calibration in
\autoref{sp:sec:results} --- the untargeted baseline that
target-informed optimization exists to beat.

Third, deployment makes data-stack sensitivity measurable.
Rerunning the companion decomposition with identical provision
dictionaries on the retired predecessor data stack moves the headline
aggregate by \$66 billion and flips the sign of the deep-poverty
response, while headline signs --- aggregate direction, a
bottom-decile loss, a rising Gini alongside a falling top-1\% share
--- agree across stacks. Incidence conclusions inherit the data
pipeline; publishing the pipeline, its calibration diagnostics, and
cross-stack comparisons alongside the policy numbers is what makes
that inheritance auditable.
